\subsection{\texorpdfstring{一元多项式的整除性 \(^{1}\)}{一元多项式的整除性 [1]}}

命题11。——设 K 是一个交换域。

\begin{enumerate}
\def\labelenumi{(\roman{enumi})}
\item
  对于每个非零理想 \(\mathfrak{a}\) （在 \(\mathbf{K}[\mathbf{X}]\) 中），存在且仅存在一个首一多项式 \(f\) （在 \(\mathbf{K}[\mathbf{X}]\) 中）使得 \(\mathfrak{a} = (f)\) 。
\item
  设 \(f_{1}\) 与 \(f_{2}\) 属于 \(\mathbf{K}[\mathbf{X}]\) ；要使 \((f_{1}) = (f_{2})\) 成立，必要且充分条件是存在非零元素 \(\lambda\) （ \(\mathbf{K}\) 中的）使得 \(f_{2} = \lambda f_{1}\) 。
\end{enumerate}

我们来证明 (ii)，所述条件的充分性是显然的。 \(f_{1}\) 与 \(f_{2}\) 生成零理想的情形是平凡的。因此设非零多项式 \(f_{1}\) 与 \(f_{2}\) 生成 K{[}X{]} 的同一个理想。那么存在多项式 \(u_{1}\) 与 \(u_{2}\) 使得 \(f_{1}=u_{1}f_{2}\) 且 \(f_{2}=u_{2}f_{1}\) ；由此可得 \(u_{1}u_{2} = 1\) ，从而 \(\deg u_{1} + \deg u_{2} = 0\) ，因此 \(\deg u_{2} = 0\) 。这样我们就证明了 \(u_{2}\) 是 \(\mathbf{K}\) 的一个非零元素。

为证明 (i)，设 \(f\) 是 \(\mathfrak{a}\) 中一个次数尽可能小的首一多项式。给定 \(g\) （在 \(\mathfrak{a}\) 中），设 \(u\) 与 \(v\) 为 \(g\) 除以 \(f\) 的欧几里得除法的商与余式；则 \(v = g - uf\) 属于 \(\mathfrak{a}\) ，且我们有 \(\deg v < \deg f\) ；若 \(v\) 非零，则存在非零元素 \(\lambda\) （ \(K\) 中的）使得 \(\lambda v\) 是首一的，而由于 \(\lambda v \in \mathfrak{a}\) ，这将与 \(f\) 的定义矛盾。于是我们有 \(\mathfrak{a} = (f)\) ；首一多项式 \(f\) （使得 \(\mathfrak{a} = (f)\) 者）的唯一性现在由 (ii) 得出。

命题12。——设 K 是一个交换域， f 和 g 是 K{[}X{]} 的两个元素。对于 K{[}X{]} 中的每个多项式 d ，以下性质是等价的：

\begin{enumerate}
\def\labelenumi{(\roman{enumi})}
\item
  多项式 \(d\) 整除 \(f\) 与 \(g\) ，且每个同时整除 \(f\) 与 \(g\) 的多项式都整除 \(d\) 。
\item
  多项式 \(d\) 整除 \(f\) 与 \(g\) ，且存在两个多项式 \(u\) 与 \(v\) 使得 \(d = uf + vg\) 。
\item
  理想之间的关系 \((d) = (f) + (g)\) 在 \(\mathbf{K}[\mathbf{X}]\) 中成立。
\end{enumerate}

多项式 d 由这些性质确定，且可相差 K 中一个非零元素的乘法。若 f 和 g 不全为零，则 \(d \neq 0\) 且 d 的次数大于等于整除 f 和 g 的每一个多项式的次数。

当 \(f\) 与 \(g\) 均为零时，性质 (i) 至 (iii) 中的每一条仅对 \(d = 0\) 成立，因此它们此时等价。以下我们假设 \(f, g\) 不全为 0 ，并用 \(\mathfrak{a}\) 表示理想 \((f) + (g)\) （ \(\mathbf{K}[\mathbf{X}]\) 中的）。

我们指出，对任意多项式 \(u\) 与 \(v\) （在 \(\mathbf{K}[\mathbf{X}]\) 中），性质 \((u) \supset (v)\) 与 \(\langle u \text{ 整除 } v \rangle\) 等价。于是断言 (ii) 等价于 \(\langle (d) \supset (f) \text{ 且 } (d) \supset (g) \text{ 且 } d \in (f) + (g) \rangle\) ，即 (iii)。显然 (ii) 蕴含 (i)。最后假设 (i) 成立；我们有 \((d) \supset (f)\) 与 \((d) \supset (g)\) ，从而 \((d) \supset \mathfrak{a}\) ；另一方面，由命题 11（IV，第 12 页），存在多项式 \(d_1\) 使得 \(\mathfrak{a} = (d_1)\) ；由于 \(d_1\) 同时整除 \(f\) 与 \(g\) ，由假设它整除 \(d\) ，从而 \((d) \subset \mathfrak{a}\) ，最后我们有 \((d) = \mathfrak{a}\) ，即 (iii)。

命题 12 的其余断言是把命题 11 应用于理想 \(\mathfrak{a} = (f) + (g)\) 而得到的直接推论。

定义 1. --- 采用命题 12 的记号，我们称 d 为 f 与 g 的一个最大公因子（简称 gcd）。若 1 是 f 与 g 的一个最大公因子，则称 f 与 g 互素，或称 f 与 g 互质。

称 f 与 g 互素，即意味着存在 K{[}X{]} 中的多项式 u 和 v 使得 \(uf + vg = 1\) 。

推论1. — 设 \(d\) 是 \(f\) 与 \(g\) 的一个最大公因子，且 \(\mathbf{K}'\) 是一个包含 \(\mathbf{K}\) 作为子域的交换域。则 \(d\) 是 \(f\) 与 \(g\) （视为 \(\mathbf{K}'[\mathbf{X}]\) 的元素）的一个最大公因子。

这由命题 12（iii）得出。

推论2. --- 设 d 是 f 与 g 的一个最大公因子。

\begin{enumerate}
\def\labelenumi{(\roman{enumi})}
\item
  如果 \(u \in \mathbf{K}[\mathbf{X}]\) ，则 du 是 fu 与 gu 的一个最大公因子。
\item
  如果 \(v \in \mathbf{K}[\mathbf{X}]\) 是（ \(\neq 0\) ） \(f\) 与 \(g\) 的一个公因子，则 \(d / v\) 是 \(f / v\) 与 \(g / v\) 的一个最大公因子。
\end{enumerate}

这由命题 12（ii）推出。

推论3. —— 设 w 是 f 和 g 的一个公因子。 w 成为 f 和 g 的最大公因子的充分必要条件是 f/w 与 g/w 互素。这由推论 2 得出。

推论4. --- 设 \(f, g, h \in \mathbf{K}[\mathbf{X}]\) 。若 \(f\) 整除 \(gh\) 且与 \(g\) 互素，则 \(f\) 整除 \(h\) 。

因为 \(f\) 整除 \(gh\) 与 \(fh\) ，因此 \(f\) 整除 \(gh\) 与 \(fh\) 的每个最大公因子，特别地整除 \(h\) （推论 2，(i)）。

推论5. --- 设 \(f, g \in \mathbf{K}[\mathbf{X}]\) 。 \(f\) 与 \(g\) 互素的充分必要条件是 \(g\) 在 \(\mathbf{K}[\mathbf{X}] / (f)\) 中的典范像可逆。

因为这一条件意味着存在 \(u, v \in \mathbf{K}[\mathbf{X}]\) 使得 \(uf + vg = 1\) 。

推论6. --- 设 \(f, g_1, g_2, \ldots, g_n \in \mathbf{K}[\mathbf{X}]\) 。若 \(f\) 与 \(g_1, g_2, \ldots, g_n\) 互素，则 \(f\) 与 \(g_1g_2\ldots g_n\) 互素。

* 推论7. --- \(f\) 与 \(g\) 互素的充分必要条件是，它们在 \(K\) 的任何扩张中都没有公共根。

因为若 \(d\) 是 \(f, g\) 的一个最大公因子，则 \(f\) 与 \(g\) 在扩张 \(\mathbf{K}'\) （ \(\mathbf{K}\) 的扩张）中的公共根就是 \(d\) 在 \(\mathbf{K}'\) 中的根。于是推论由 \(V\) ，第 21 页，命题 4 得出。*

